% PCB 배치 부록 (그림은 report/figures/make_pcb_figures.py로 생성).
\chapter{PCB 배치}\label{app:pcb}

회로도 부록(\ref{app:schematics})의 다섯 보드를 기판에 배치한 결과다. 그림은 앞면의 배치도로, 외곽, 부품의 외형,
실크, 부품이 차지하는 경계를 그렸다. \emph{배선은 아직 하지 않았다.} 보드 파일에는 외곽, 장착 구멍, 모든 부품의
자리와 방향, 회로도에서 가져온 넷(어느 패드끼리 이어야 하는가), 설계 규칙이 들어 있고 동선은 하나도 없다.

\section{배치를 정한 원칙}

\paragraph{왜 배치까지만 했는가.} 이 보드들의 배선은 같은 모양의 되풀이다. 시프터는 릴레이 130개가, 테스트
레지스터는 레인 26개가, 계수기는 자리 다섯이 같은 회로다. 되풀이되는 단위 하나의 배선을 사람이 정하고 나머지에
옮기는 것이 자동 배선보다 읽기 쉽고 고치기 쉬운 기판을 만든다. 배치는 그 전제다. 되풀이 단위가 기판 위에서도
같은 모양으로 놓여 있어야 한 단위의 배선을 옮길 수 있다. 그래서 배치도 회로도처럼 스크립트가 만든다
(\code{gen\_*\_pcb.py}). 부품의 자리는 ``몇째 단의 몇째 선''에서 계산되는 좌표이고, 같은 입력은 바이트 단위로 같은
파일을 낸다.

\paragraph{공통 규칙.}
\begin{itemize}
\item \textbf{신호가 흐르는 방향으로 놓는다.} 입력 커넥터에서 출력 커넥터까지 선이 되돌아가지 않게 한다. 26선이
나란히 가는 보드에서는 되돌아가는 선 하나가 26번의 교차가 된다.
\item \textbf{커넥터는 가장자리에, 같은 종류는 같은 변에.} 케이블이 보드 위를 가로지르지 않는다. 릴레이식의 세
보드는 입력이 왼쪽, 출력이 오른쪽이어서 카드 케이지에 나란히 꽂았을 때 케이블이 한 방향으로 흐른다.
\item \textbf{코일의 다이오드는 코일 옆에.} 코일이 꺼질 때의 전류가 도는 고리를 작게 한다.
\item \textbf{릴레이식의 세 보드는 같은 카드 크기}(250 × 228~mm, 구멍은 모서리에서 5~mm, M3 넷)다. 시프터 보드가
이 크기를 정했고 계수기 두 장이 따랐다. 케이지의 슬롯이 한 종류가 된다.
\end{itemize}

\paragraph{확인하는 것.} 배치만 한 기판에서도 틀릴 수 있는 것이 있다. 부품이 빠지거나, 풋프린트가 회로도와
다르거나, 패드에 붙은 넷이 회로도와 다르거나, 부품끼리 겹치는 것이다. 보드마다 검증 스크립트가 다음을 본다.
(1) 회로도의 부품이 모두 기판에 있고 풋프린트가 같다. (2) 패드마다의 넷이 회로도의 넷리스트와 정확히 같다.
(3) KiCad의 설계 규칙 검사(DRC)에서 위반이 0건이고, 남는 것은 ``아직 이어지지 않은 연결''뿐이다. (4) 패드 둘의
넷을 바꾼 기판을 넣으면 검사가 실패한다. (5) 두 번 생성한 파일이 같다.

\begin{table}[ht]
\centering
\caption{배치한 보드. 남은 연결은 배선해야 할 패드 쌍의 수다.}
\label{tab:pcb-boards}
\small
\begin{tabularx}{\linewidth}{>{\raggedright\arraybackslash}Xlrrr}
\toprule
보드 & 크기, 층 & 부품 & 넷 & 남은 연결 \\
\midrule
릴레이 시프터 & 250 × 228 mm, 4층 & 275 & --- & 1,120 \\
릴레이 계수기(공통) & 250 × 228 mm, 4층 & 261 & 211 & 539 \\
릴레이 계수기(유닛) & 250 × 228 mm, 4층 & 196 & 203 & 445 \\
테스트 레지스터 & 250 × 140 mm, 2층 & 258 & 288 & --- \\
단위 시험기 & 290 × 116 mm, 4층 & 446 & 445 & --- \\
\bottomrule
\end{tabularx}
\end{table}

(부품 수는 장착 구멍을 뺀 것이다. 단위 시험기만 구멍 여섯을 포함한다. 빈칸은 세지 않은 값이다. DRC의 보고서는
연결 안 된 항목을 499건까지만 적으므로, 남은 연결은 기판의 연결 정보에서 따로 센다.)

\newcommand{\pcbfigure}[3]{%
  \IfFileExists{figures/generated/pcb_#1.pdf}{%
    \begin{figure}[p]
      \centering
      \includegraphics[width=\linewidth,height=0.8\textheight,keepaspectratio]{figures/generated/pcb_#1.pdf}
      \caption{#2}\label{fig:pcb-#1}
    \end{figure}}{}%
  #3}

% ======================================================================
\section{릴레이 시프터 보드}

\paragraph{배치.} 릴레이를 5열 × 26행의 격자로 놓았다(그림~\ref{fig:pcb-shifter-board}). 열이 단(왼쪽부터 1, 2, 4,
8, 16칸)이고 행이 선(위가 A)이다. 릴레이를 90° 돌려 한 줄의 핀(코일, 놓임 쪽, 공통, 동작 쪽)이 가는 길의 행이
되고 5.08~mm 위의 다른 줄이 오는 길의 행이 되게 했다. 그러면 신호가 왼쪽 커넥터에서 들어와 다섯 열을 차례로
지나 오른쪽 커넥터로 나간다. 입력 둘과 제어 커넥터가 왼쪽 변, 출력 둘이 오른쪽 변이다. 다이오드는 릴레이마다
코일 핀의 왼쪽에 세로로, 단의 표시등은 열의 아래에 있다.

\paragraph{크기.} 처음에는 160 × 200~mm를 생각했으나 들어가지 않는다. 릴레이 하나가 차지하는 경계가
7.9 × 15.3~mm여서 한 열 26개만으로 208~mm다. 높이 228~mm는 여기서 나온다. 폭은 배선이 정한다. 한 단은 26선을
$s$칸 아래로(오는 길은 위로) 옮겨 다음 열에 넘기는데, 이 옮김이 열 사이에서 세로로 달리는 선이 된다. 이동량이
클수록 겹치는 세로 선이 많아 16칸 단에서는 한쪽에 약 21가닥이다. 열 사이를 38~mm 비워 이 선들의 통로로 두었다.
배선을 해 보고 통로를 줄일 수 있으면 폭이 준다.

\paragraph{층과 규칙.} 4층이다. 가는 길은 앞면(가로)과 안쪽 첫째 층(세로), 오는 길은 뒷면(가로)과 안쪽 둘째
층(세로)에 두도록 계획했다. 두 길이 층으로 나뉘므로 서로 건너지 않는다. 규칙은 흔한 4층 제작 사양에
맞췄다: 선폭과 간격 0.15~mm, 비아 0.5/0.25~mm. 신호선은 0.2~mm(선 하나에 약 10~mA), 코일의 제어선과 0~V는
0.6~mm 이상(한 단에 최대 약 0.25~A, 보드 전체의 복귀 약 0.9~A)이다.

\paragraph{확인.} 부품 275개와 패드 1,482개의 넷이 회로도와 같다. DRC 위반 0건, 남은 연결 1,120개.

\pcbfigure{shifter-board}{릴레이 시프터 보드의 배치(250 × 228\,mm, 4층). 릴레이 G6S-2를 5단 × 26선 격자로 놓아
신호가 왼쪽 입력 커넥터에서 오른쪽 출력 커넥터로 흐른다. 역기전력 다이오드는 각 코일 옆, 단 표시등은 각 열 아래.
열 사이의 빈 띠(38\,mm)는 26선을 옮겨 잇는 세로 배선의 통로다.}{}

% ======================================================================
\section{릴레이 계수기 보드 (공통, 유닛)}

\paragraph{배치.} 26진 한 자리(릴레이 11개, 다이오드 64개, 표시등 5개)를 98 × 84~mm의 블록 하나로 만들고 이
블록을 되풀이했다(그림~\ref{fig:pcb-counter-common}, \ref{fig:pcb-counter-unit}). 블록 안은 비트마다 한 열(피치
16~mm)이다. 위에서부터 종 릴레이, 주 릴레이, 그 비트의 다이오드(두 줄), 표시등이다. 올림 사슬이 비트 0에서 4로
가므로 사슬이 블록을 왼쪽에서 오른쪽으로 한 번 지나간다. 여섯째 열에 25를 알아보는 릴레이가 있다. 블록의
어느 자리를 보면 어느 비트의 무엇인지 읽힌다.

공통 보드에는 블록 셋($l$, $m$이 위, $a_3$이 아래 왼쪽)과 자리 사이의 규칙 다이오드가 있고 커넥터 둘(클럭,
유닛 버스)이 왼쪽 변에 있다. 유닛 보드에는 블록 둘($a_2$, $a_1$)과 중계 릴레이 10개가 있고, 버스와 다이얼
커넥터가 왼쪽, 시프터 넉 장으로 가는 제어 커넥터 넷이 오른쪽 변에 있다.

\paragraph{크기.} 두 장 모두 시프터 보드와 같은 250 × 228~mm다. 공통 보드는 블록 셋이 판의 절반쯤을 차지해
더 작게 만들 수도 있었다(약 250 × 190~mm). 카드 크기를 하나로 맞추는 쪽과 배선의 여유를 택했다.

\paragraph{층과 규칙.} 4층, 제작 규칙은 시프터 보드와 같다. 넷을 두 무리로 나눴다. 신호(0.2~mm)와 레일
(0.6~mm, 비아 0.8/0.4~mm)이다. 계수기에서는 클럭의 두 상, 올림 레일, 코드선이 모두 릴레이 코일을 먹이는
전원이므로 레일에 넣었다. 가장 깊은 길은 다이오드 넷을 지나 코일 전압의 여유가 작다. 레일의 저항이 그 여유를
더 깎지 않아야 한다.

\paragraph{확인.} 공통 보드는 부품 261개, 패드 750개, 유닛 보드는 196개, 648개의 넷이 회로도와 같다. DRC 위반
0건. 남은 연결은 539개와 445개다. 패드 둘의 넷을 바꾼 것과 부품을 뒷면으로 뒤집은 것을 검사가 잡았다. 두 번
생성한 파일이 같다.

\pcbfigure{counter-common}{릴레이 계수기 공통 보드의 배치(250 × 228\,mm, 4층, 시프터 보드와 같은 카드 크기).
26진 한 자리(주종 릴레이 5쌍 + 25 감지 릴레이 + 다이오드)가 98 × 84\,mm 블록 하나이고, 이 블록이 $l$, $m$, $a_3$
세 자리로 되풀이된다. 자리 사이의 규칙 다이오드와 커넥터는 왼쪽 가장자리.}{}

\pcbfigure{counter-unit}{릴레이 계수기 유닛 보드의 배치(250 × 228\,mm, 4층). 다이얼로 설정하는 $a_2$, $a_1$ 두 자리
블록과 $a_3$/$a_4$ 버퍼 릴레이 10개. 오른쪽 가장자리의 D-sub 9 네 개가 시프터 보드 네 장의 제어 입력으로 간다.}{}

% ======================================================================
\section{테스트 레지스터 보드}

\paragraph{배치.} 레인 26개를 3행 × 9열의 격자로 놓았다(그림~\ref{fig:pcb-test-register}). 레인 하나가 글자
하나이고 모양이 모두 같다: 릴레이 옆에 차단 다이오드와 코일 다이오드, 게이트 분압 저항 둘, 표시등과 저항,
MOSFET. 위쪽 띠에는 한 개씩뿐인 것들이 있다. 타이밍·통과·정지·해제 릴레이 여덟과 그 다이오드, 오른쪽 위에
전원과 시험 선과 패널의 단자, 퓨즈, 커패시터다. D-sub 세 개(26선 입력, 뱅크 제어, 표시)는 아래 변에 나란히
있어 케이블이 한쪽으로만 나온다.

레인을 격자로 놓은 이유는 레인들을 잇는 넷이 많기 때문이다. 유지 버스, 0~V, 표시 쪽 0~V는 모든 레인에
가고, 접점 사슬은 레인을 A에서 Z까지 차례로 꿴다. 격자에서는 사슬이 행을 따라 갔다가 다음 행으로 넘어가는
뱀 모양이 되고, 버스는 행마다 한 줄이 된다.

\paragraph{크기와 층.} 250 × 140~mm, 2층이다. 시프터와 달리 26선이 서로를 넘어 옮겨지는 배선이 없다. 레인
안의 연결은 짧고 레인 사이는 버스와 사슬뿐이어서 두 층으로 된다. 전류가 1~A에 가까운 넷(0~V, 유지 버스,
$+24$~V, 표시 쪽)은 굵게 배선한다.

\paragraph{확인.} 부품 258개, 넷 288개가 회로도와 같다. DRC는 회로도와의 일치 검사를 포함해 위반 0건이다.

\pcbfigure{test-register}{테스트 레지스터 보드의 배치(250 × 140\,mm, 2층). 감지 릴레이 26개와 그 다이오드·램프
구동이 3 × 9 격자로 같은 모양을 되풀이하고, 통과·정지·타이밍 릴레이와 전원·패널 단자가 위쪽, D-sub 세 개가
아래 가장자리에 있다.}{}

% ======================================================================
\section{스크램블러 단위 시험기}

\paragraph{배치.} 이 보드의 배치는 절연이 정한다(그림~\ref{fig:pcb-unit-tester}). 기판을 가로지르는 선 하나를
경계로 삼아 위쪽을 USB 영역, 아래쪽을 24~V 현장 영역으로 나눴다. 두 영역을 건너는 부품 22개(디지털 절연기
둘, 고속 광커플러 하나, 4채널 광커플러 열아홉)가 그 선 위에 한 줄로 걸터앉는다. 경계의 띠에는 다른 구리가
없다. 건너는 길이 이 부품들뿐이라는 것이 눈으로 보인다.

위쪽에는 Pico(USB 단자가 왼쪽 변)와 감지용 시프트 레지스터 열 개가 한 줄로, 광커플러마다 그 출력의 저항이
바로 위에 있다. 아래쪽은 노드 여덟 개의 묶음이 세로 한 칸씩 일곱 칸이다. 한 칸 안에서 위에서 아래로 감지 저항과
보호 다이오드, 1~k$\Omega$ 여덟 개, 구동기와 구동용 시프트 레지스터 순서다. 감지가 경계 가까이, 구동이 커넥터
가까이 있다. 아래 변에 D-sub 37 세 개가 있고 노드 0--25가 A 커넥터 쪽에 가깝다. 왼쪽 아래에 전원이 있다.

\paragraph{절연 거리.} 두 영역 사이의 연면 거리와 공간 거리를 4.5~mm 이상으로 잡고 DRC의 규칙으로 걸었다.
현장 영역의 넷과 USB 영역의 넷 사이가 이보다 가까우면 위반이 된다. 규칙이 실제로 일하는지는 값을 12~mm로
올려 위반이 나오는 것으로 확인했다. 실제의 최소는 광커플러의 패드 사이 4.9~mm다. 경계선의 양 끝, 부품이 없는
구간에는 폭 1~mm의 슬롯을 내어 기판 표면을 따라가는 길을 끊었다. 4.5~mm는 24~V 회로의 기능 절연에 여유를 둔
설계 목표이고, 내전압 시험을 한 값은 아니다.

\paragraph{크기와 층.} 290 × 116~mm, 4층(신호, 0~V, 전원, 신호)이다. 폭은 경계에 한 줄로 놓인 절연 부품 22개
(피치 11.4~mm)와 아래 변의 D-sub 세 개(피치 75~mm)가 정한다. 0~V 층은 두 영역으로 갈라 채운다.

\paragraph{확인.} 부품 446개, 넷 445개가 회로도와 같다. DRC 위반 0건.

\pcbfigure{unit-tester}{스크램블러 단위 시험기 보드의 배치(290 × 116\,mm, 4층). 왼쪽 위가 Pico와 USB(저전압
쪽), 가운데 가로줄이 절연 장벽(디지털 절연기와 광커플러, 양 끝 슬롯), 그 아래가 24\,V 구동·감지 노드 7묶음과
전원, 아래 가장자리에 D-sub 37 세 개. 장벽을 건너는 거리는 4.5\,mm 이상이 되도록 DRC 규칙으로 강제했다.}{}

% ======================================================================
\section{남은 일}

\begin{itemize}
\item \textbf{인터페이스 결정의 반영.} 배치를 마친 뒤에 커넥터와 접지의 규약을 정했다
(\code{docs/decisions/0003-interfaces.md}). 보드마다 고칠 곳이 있다(부록~\ref{app:schematics}의 각 절 끝). 단위
시험기의 커넥터를 소켓에서 핀으로, 계수기에 전원 단자, 테스트 레지스터의 패널 단자와 입력 다이오드 같은 것이다.
생성기를 고쳐 다시 만들면 배치와 검증이 따라온다.
\item \textbf{배선.} 사람이 한다. 되풀이 단위 하나를 배선해 옮기는 순서를 권한다. 배선과 구리면을 채운 뒤에
DRC를 다시 돌린다.
\item \textbf{그림.} 배선한 뒤 저장소 루트에서 \code{report/figures/make\_pcb\_figures.py}를 다시 실행하면 이 부록의
그림이 바뀐다.
\item \textbf{부품의 확인.} 릴레이의 코일 극성과 핀 기능, 퓨즈와 절연기의 품번은 데이터시트로 확인하지 못한 것이
있다. 주문 전에 확인한다.
\end{itemize}

\pcbfigure{contact-board}{기어식 접점 보드의 배치(110 × 110\,mm, 4층). Omniball 접점 104개가 드럼 PCB의 네 패드 링과 마주 보도록 놓이고(접점 쪽에서 반시계 방향), 네 모서리의 2×13 헤더가 위아래 자리와 리본으로 이어진다.}{}
\pcbfigure{diagonal-d2}{대각선 보드 D2의 배치(369 × 130\,mm). 글자마다 D-sub 37 잭 두 개가 110 × 30\,mm 격자로 놓이고, 양쪽 가장자리의 64핀 헤더 세 개가 이웃 보드로 간다. 다섯 장을 이으면 1,562\,mm가 된다(계획 1,430\,mm보다 길다).}{}
\pcbfigure{patch-panel}{메뉴 패치 패널의 배치(339.5 × 160\,mm). 공통 잭 16개, 레지스터 잭, 탐색 잭과 시험 플러그 구멍 26개.}{}
